\section{Standard weighted dependence measures}
\label{sec:standard-measures}

This section records four familiar coefficients in the case-weight notation
of Section~\ref{sec:framework}. The formulas are included to fix conventions
and to provide comparators for the less standard developments in later
sections. Every coefficient is invariant to a common positive rescaling of the
original weights, and zero-mass observations make no contribution.

\subsection{Pearson correlation and Spearman's rho}

The weighted means, variances, and covariance are

\begin{align*}
  \bar X_w
  &= \sum_{i=1}^n p_iX_i,
  &
  \bar Y_w
  &= \sum_{i=1}^n p_iY_i,\\
  s_{X,w}^2
  &= \sum_{i=1}^n p_i(X_i-\bar X_w)^2,
  &
  s_{Y,w}^2
  &= \sum_{i=1}^n p_i(Y_i-\bar Y_w)^2,\\
  s_{XY,w}
  &= \sum_{i=1}^n p_i(X_i-\bar X_w)(Y_i-\bar Y_w).
\end{align*}

The weighted Pearson coefficient is

\begin{align*}
  \rho_{P,w}
  &= \frac{s_{XY,w}}{s_{X,w}s_{Y,w}}.
\end{align*}

It is undefined if either weighted marginal variance is zero. Apart from the
replacement of empirical averages by weighted empirical averages, this is the
classical product-moment coefficient associated with Pearson's work on
correlation and regression \citep{pearson1895}.

Weighted Spearman correlation is defined as weighted Pearson correlation
between the average scores from Section~\ref{sec:framework}:

\begin{align*}
  \rho_{S,w}
  &= \frac{
      \sum_{i=1}^n p_i(R_i^X-\bar R_w^X)(R_i^Y-\bar R_w^Y)
    }{
      \sqrt{\sum_{i=1}^n p_i(R_i^X-\bar R_w^X)^2}
      \sqrt{\sum_{i=1}^n p_i(R_i^Y-\bar R_w^Y)^2}
    },\\
  \bar R_w^X
  &= \sum_{i=1}^n p_iR_i^X,
  &
  \bar R_w^Y
  &= \sum_{i=1}^n p_iR_i^Y.
\end{align*}

Uniform weights recover the ordinary coefficient with midranks for ties
\citep{spearman1904}. Sorting dominates the computation, so the weighted
coefficient can be obtained in $O(n\log n)$ time and $O(n)$ memory.

\subsection{Kendall's tau}

Kendall's coefficient is based on pairwise concordance
\citep{kendall1938}. With $\operatorname{sgn}(0)=0$, define the weighted
concordance difference and the masses of pairs untied in each margin as

\begin{align*}
  S_w
  &= \sum_{1\leq i<j\leq n}p_ip_j
     \operatorname{sgn}(X_i-X_j)\operatorname{sgn}(Y_i-Y_j),\\
  Q_{X,w}
  &= \sum_{1\leq i<j\leq n}p_ip_j\mathbf{1}(X_i\neq X_j),\\
  Q_{Y,w}
  &= \sum_{1\leq i<j\leq n}p_ip_j\mathbf{1}(Y_i\neq Y_j).
\end{align*}

The weighted tau-b coefficient is

\begin{align*}
  \tau_w
  &= \frac{S_w}{\sqrt{Q_{X,w}Q_{Y,w}}}.
\end{align*}

Thus each pair receives the product of its two case masses. The denominator
removes pairs tied in either margin and reduces to $e_2(p_1,\ldots,p_n)$ in
both factors when there are no ties. Uniform weights give the ordinary sample
tau-b. Multiplicative pair weights are also used in other weighted Kendall
constructions \citep{shieh1998}, although those constructions need not share
the present case-weight interpretation.

Direct evaluation is quadratic. As in the unweighted algorithm of
\citet{knight1966}, discordant pairs can instead be counted during merge sort.
Carrying the case masses through the merge replaces an inversion count by a
sum of pair products and preserves $O(n\log n)$ running time.

\subsection{Blomqvist's beta}

Blomqvist's beta is a quadrant-based coefficient \citep{blomqvist1950}. It
records whether observations fall in matching quadrants about the marginal
centers. To specify the finite-sample center used here, sort the observations
by $X$ and let $R_{(k)}^X$ denote the weighted
average score attached to the $k$th ordered value. With

\begin{align*}
  k_X
  &= \min\left\{k:R_{(k)}^X\geq e_2(p_1,\ldots,p_n)\right\},
\end{align*}

define

\begin{align*}
  m_{X,w}
  &=
  \begin{cases}
    X_{(k_X)},
      & R_{(k_X)}^X=e_2(p_1,\ldots,p_n),\\
    \left(X_{(k_X-1)}+X_{(k_X)}\right)/2,
      & R_{(k_X)}^X>e_2(p_1,\ldots,p_n).
  \end{cases}
\end{align*}

Define $m_{Y,w}$ analogously. This pair-mass center agrees with the ordinary
sample median under uniform weights and makes the treatment of unequal case
weights explicit. It is not interchangeable with every convention called a
weighted median.

The sample coefficient is

\begin{align*}
  \beta_w
  &= 2\sum_{i=1}^n p_i
     \left\{
       \mathbf{1}(X_i\leq m_{X,w},\,Y_i\leq m_{Y,w})
       +
       \mathbf{1}(X_i>m_{X,w},\,Y_i>m_{Y,w})
     \right\}-1.
\end{align*}

The lower-left quadrant includes observations on either center and the
upper-right quadrant is strict. For continuous population marginals this
corresponds to the familiar copula functional
$4C(1/2,1/2)-1$. Computing both centers by sorting takes
$O(n\log n)$ time.

\subsection{Summary}

\begin{table}[ht]
  \centering
  \caption{Standard case-weighted dependence measures.}
  \label{tab:standard-measures}
  \begin{tabular}{llll}
    \toprule
    Measure & Basic object & Tie convention & Time \\
    \midrule
    Pearson & weighted moments & not applicable & $O(n)$ \\
    Spearman & weighted average ranks & pair-mass average & $O(n\log n)$ \\
    Kendall & weighted concordant pairs & tau-b denominator & $O(n\log n)$ \\
    Blomqvist & matching median quadrants & pair-mass center & $O(n\log n)$ \\
    \bottomrule
  \end{tabular}
\end{table}

Pearson, Spearman, Kendall, and Blomqvist coefficients are symmetric in the two
variables. The first three are undefined when their displayed denominator
vanishes. Their independence-test transformations and the more elaborate tie
adjustment for Kendall's statistic are deferred to Section~5.
